\cleardoublepage
\chapternonum{摘要}

空地协同系统结合了无人机的空中视角与高机动性，以及无人车的地面负载与持续作业能力，在搜索救援、环境探索等任务中具有广阔应用前景。面向室内、城市峡谷、林下等卫星信号受遮挡或不可用的场景，传统依赖 GNSS 的导航难以直接适用；即便采用视觉、激光等机载自主定位方案，也常因视角变化、纹理退化与快速机动出现累计漂移，或带来较高的传感与计算能耗开销。相比之下，地面平台具备更强载荷与续航能力，可搭载更高性能的传感器与计算单元，更容易获得稳定的里程计与地图约束，从而为协同定位提供更可靠的参考框架。

基于这种互补性，空地系统通常通过空地间相对观测与信息共享实现联合定位与导航：地面平台提供稳定参考与环境感知，无人机负责完成任务，并利用相对测距、互观约束等抑制漂移。然而协同依赖持续的数据交互与良好的观测几何，当障碍物造成视线遮挡时，通信与相对观测同时退化，进而影响定位精度与协同规划可靠性。因此，如何在未知复杂环境下实现高精度、低漂移的空地协同定位，并在规划中主动维持可见性、生成安全高效的协同轨迹，仍是工程落地的关键问题。为此，本文面向卫星拒止环境，研究空地异构集群的相对定位架构与可见性驱动的协同轨迹规划方法。本文的主要工作如下：

（1）提出了一种基于动态测距基站的空地相对定位框架。针对传统测距定位依赖固定基站、部署成本高且适应性差的问题，本文将测距基站部署在移动无人车上，并以无人车构建的激光建图全局坐标系作为统一参考。通过高频测距时序协议提升测距更新率，并采用误差状态扩展卡尔曼滤波将多边定位解算结果与机载惯性测量数据融合，实现了无人机在无卫星条件下的高频、稳定相对位姿估计。此外，构建了基于三维栅格占据图的共享地图机制，为空地协同规划提供统一的地图输入。

（2）提出了一种基于点云聚类的面向视线维持的空地协同规划方法。针对障碍物遮挡导致的测距（通信）链路退化与碰撞风险，本文首先对环境点云进行聚类，并以最小包络椭圆柱对障碍物进行紧凑凸几何近似；在此基础上构造可微的几何避障约束，并设计引导目标、空地视线连通性与空地编队代价，然后统一纳入非线性模型预测控制框架。仿真结果表明，该方法在未知简单环境中能够实时生成平滑、安全且兼顾视线连通性的协同轨迹。

（3）提出了一种基于序列可视空间的两阶段空地协同规划方法。为提升稠密障碍环境下的求解稳定性与鲁棒性，本文构建了“全局一致性规划—局部安全执行”的两阶段架构：全局层基于动力学可行的路径搜索生成引导序列，并在路径上构建时序化的可视空间，将复杂的避障与视距约束转化为可微的可见性代价，从而生成满足可见性与队形要求的全局参考路点；局部层采用多项式轨迹参数化，在跟踪全局参考点的同时，结合局部避障力场，实现动力学可行且安全的连续轨迹输出。仿真与实物实验验证了该方法在复杂环境下具有更高的任务成功率与鲁棒性。

\vspace{1em}
\noindent \textbf{关键词：} 空地协同；卫星拒止环境；相对定位；可见性规划


\cleardoublepage
\chapternonum{Abstract}

Air-ground collaborative systems leverage the superior mobility of Unmanned Aerial Vehicles (UAVs) and the significant endurance and payload capacity of Unmanned Ground Vehicles (UGVs), offering broad prospects in search, rescue, and exploration. In GNSS-denied scenarios such as indoors or under forest canopies, traditional satellite-dependent navigation is inapplicable, and onboard-only sensing often suffers from cumulative drift and high energy overhead. In contrast, ground platforms carry high-performance sensors to provide stable odometry and mapping constraints, serving as a reliable reference frame for the collaborative system.

Based on this complementarity, air-ground system achieve joint navigation through relative observations and information sharing. However, such collaboration relies on continuous line-of-sight (LOS) connectivity. Obstacle occlusion degrades both communication and ranging links, undermining localization accuracy and planning reliability. Therefore, achieving high-precision relative localization and maintaining visibility through safe, efficient trajectory generation remain critical challenges. This thesis investigates a relative localization architecture and visibility-driven collaborative planning methods for heterogeneous air-ground clusters in GNSS-denied environments. The main contributions are as follows:

(1) A relative localization framework based on dynamic ranging base stations is proposed. Addressing the high deployment costs and poor adaptability of traditional fixed-station ranging systems, this work deploys ranging base stations on mobile UGVs and utilizes the global coordinate system established by UGV-based LiDAR mapping as a unified reference. Ranging update rates are enhanced through a high-frequency ranging timing protocol. An Error-State Extended Kalman Filter (ESKF) is employed to fuse multilateration results with airborne inertial measurement data, achieving high-frequency and stable relative pose estimation for UAVs under GNSS-denied conditions. Furthermore, a shared mapping mechanism based on 3D occupancy grid maps is constructed to provide a unified map input for air-ground collaborative planning.

(2) A visibility-aware air-ground collaborative planning method based on point cloud clustering is proposed. To tackle the degradation of ranging/communication links and collision risks caused by obstacle occlusion, environmental point clouds are first clustered and approximated as compact convex geometries using minimum bounding ellipsoidal cylinders. On this basis, differentiable geometric obstacle avoidance constraints are constructed, and costs for guiding goals, LOS connectivity, and air-ground formation maintenance are designed and integrated into a Nonlinear Model Predictive Control (NMPC) framework. Simulation results demonstrate that this method can generate smooth, safe, and LOS-consistent collaborative trajectories in real-time within unknown simple environments.

(3) A two-stage air-ground collaborative planning method based on sequential visible spaces is proposed. To improve solving stability and robustness in dense obstacle environments, a "Global Consistency Planning – Local Safety Execution" two-stage architecture is developed. The global layer generates kinodynamically feasible guiding sequences through path searching and constructs temporal visible spaces along the path, converting complex obstacle avoidance and LOS constraints into differentiable visibility costs to generate global reference waypoints that satisfy visibility and formation requirements. The local layer utilizes polynomial trajectory parameterization to track global waypoints while incorporating local obstacle avoidance force fields, outputting kinodynamically feasible and safe continuous trajectories. Simulation and real-world experiments verify that this method achieves higher task success rates and robustness in complex environments.

\vspace{1em}
\noindent \textbf{Keywords:} air–ground cooperation; GNSS-denied environments; relative localization; visibility-aware planning